% !TeX root = compte-rendu.tex
% !TeX spellcheck = fr


\section{Formalisation Mathématiques}

Nous utiliserons les notations de \textit{la mesure du pouvoir de vote} \cite{mesure_pouvoir}. Notons $\V$ l'ensemble des électeurs ou votants. Nous prendrons toujours $\V = [\![0;N]\!]$ avec $N \in \mathbb{N}$.

\subsection{Jeu coopératif}

\begin{definition}
	Un \textit{jeu de vote pondéré} est un vote où chaque électeur ne peut voter que "pour" ou "contre" la proposition. C'est un couple $(q, \left(w_i\right)_{i \in \V})$ où $q$ est le \textit{quota} nécessaire pour que la proposition soit acceptée et $(\left(w_i\right)_{i \in \V})$ est le poids de chaque électeur dans ce vote.
\end{definition}

\begin{definition}
	Un \textit{coalition} est dite \textit{gagnante} si
	\begin{equation}
		\sum_{i \in S} w_i \geqslant q
	\end{equation}
	c'est-à-dire si elle peut faire accepter une proposition et dite \textit{perdante} sinon.
	Notons $\mathcal{G}$ l'ensemble des coalitions gagnantes et $\mathcal{G}_i$ l'ensemble de celles contenant l'électeur $i$. \\
\end{definition}

Posons $v : \mathcal{P}\left(\V\right) \rightarrow \left\{0,1\right\}$ la fonction indicatrice de $\mathcal{P}\left(\mathcal{G}\right)$ :
\begin{equation}
	v(S) = \left\{ \begin{array}{cl}
		1 &\text{ si } \sum_{i \in S} w_i \geqslant q \\
		0 &\text{ sinon}
	\end{array} \right.
\end{equation}
Un jeu de vote pondéré peut être représenté par $(\V, v)$ (bien que certaines fonction $v : \mathcal{P}\left(\V\right) \rightarrow \left\{0,1\right\}$ ne représente aucun jeu de vote pondéré).

\subsection{Indice de pouvoir}

Le but des indices de pouvoir est de mesurer la capacité de chaque électeur à modifier le résultat du vote.

\begin{definition}
	Un électeur est \textit{décisif} dans une coalition gagnante si lorsque cet électeur se retire de la coalition, celle-ci devient perdante. Notons la décidabilité ou l'indice de Banzhaf brut $d_i(v)$ le nombre de coalition où l'électeur $i$ est décisif.
	\begin{equation}
		d_i(v) = \sum_{S \subset \V, \ i \in S} v(S) - v(S \setminus \{i\})
	\end{equation}
\end{definition}

\begin{definition}
	L'indice de Banzhaf normalisé $B_i(v)$ est définit par
	\begin{equation}
		B_i(v) = \frac{d_i(v)}{\sum_{j=1}^{n} d_j(v)}
	\end{equation}
	L'indice de Banzhaf relatif $B_i(v)$ est définit par
	\begin{equation}
		B_i(v) = \frac{d_i(v)}{2^{|\V| - 1}}
	\end{equation}
\end{definition}

\begin{algorithm}[H]
	\caption{Calcul de l'ensemble des parties d'un ensemble}
	\Entree{Un ensemble E (sous forme de liste)}
	\Sortie{Ensemble des parties}
	\lSi{E est vide}{ \Retour{$\{\emptyset\}$} }
	\lSi{E ne contient qu'un seul élément}{ \Retour{$\{\emptyset, E\}$} }
	Choisir $e \in E$ \\
	
$F \gets Parties(\{ E \setminus \{e\} \})$ \quad \emph{Appel récursif} \\
	$P \gets \emptyset$ \\
	\Pour{$f \in F$}{
		Ajouter à $P$ les ensembles $f$ et $f \cup \{e\}$
	}
	\Retour{P}
\end{algorithm}

En notant $c_n$ la complexité de cette algorithme pour en ensemble de cardinal $n$, nous avons $c_0 = 1$, $c_1 = 2$ et $\forall n \geqslant 1, c_{n+1} = c_n + O(2^n)$. D'où $c_n - c_2 = \sum_{k=2}^{n-1} c_{k+1} - c_k$. Finalement, la complexité temporelle est en $O(n 2^n)$.

\bigbreak

\begin{algorithm}[H]
	\caption{Calcul naïf de l'indice de Banzhaf brut}
	\Entree{Un jeu de vote pondéré $(q, (w_v)_{v \in \mathcal{V}})$}
	\Sortie{Indice de Banzhaf brut de tous les électeurs}
	$G \gets \emptyset$ \quad \emph{Ensemble des parties gagnantes} \\
	\Pour{$S \in Parties(\mathcal{V})$}{
		$somme \gets \sum_{v \in S} w_v$ \\
		\lSi{$somme \leqslant q$}{ Ajouter à $G$ le tuple $(S, somme)$ }
	}
	$D \gets (0)_{v \in \mathcal{V}}$ \quad \emph{Dictionnaire associant à chaque électeur sa décidabilité} \\
	\Pour{$(S, somme) \in G$}{
		\Pour{$v \in S$}{
			\lSi{$(somme - w_v) < q$}{$D_v \gets D_v + 1$}
		}
	}
	$sommeD \gets \sum_{v \in S} D_v$ \quad \emph{Somme des décidabilités} \\
	\uSi{$sommeD = 0$}{\Retour{$(0)_{v \in \mathcal{V}}$}}
	\Sinon{\Retour{$(\nicefrac{d_v}{sommeD})_{v \in \mathcal{V}}$}}
\end{algorithm}

Dans \cite{algorithm_weight}, Dennis Leech propose

Pour l'implémentation de notre algorithme nous avons choisi d'utiliser :
\begin{equation}
	w_0 = d
\end{equation}
\begin{equation}
	\lambda(i) = \frac{0,\!83 q}{\ln(i)+1}
\end{equation}